\section{The Event Bus}

The EventBus is the nervous system of the \Core{}.  It implements a
publish-subscribe message bus with persistent storage, dead-letter queues,
and replay capability.  All other modules register as subscribers to
specific event types, ensuring full decoupling.

\subsection{API and Configuration}

The bus exposes a minimal asynchronous API:

\begin{lstlisting}[language=Python, caption={EventBus core API}]
# Emitting an event
event_id = await bus.emit(
    "chaos_state_changed",          # event_type
    {"state": "diverging"},         # data
    source="chaos_detector",        # source module
    priority=EventPriority.HIGH,
    ttl=30.0,                       # time-to-live in seconds
    metadata={"layer": "core"}
)

# Subscribing to an event type
bus.subscribe(
    "chaos_state_changed",
    callback=my_callback,
    filter_func=lambda ev: ev.data["state"] == "critical",
    retry_policy=RetryPolicy(max_retries=3, backoff_factor=2.0)
)
\end{lstlisting}

Configuration is managed via the \texttt{EventBusConfig} dataclass:

\begin{lstlisting}[language=Python, caption={EventBus configuration}]
@dataclass(frozen=True)
class EventBusConfig:
    node_id: str = "node_1"
    max_history: int = 10000
    max_queue_size: int = 10000
    batch_size: int = 10
    enable_persistence: bool = True
    persistence_path: str = "events.db"
    enable_distributed: bool = False
    distributed_backend: str = "redis"
    serializer: str = "json"
    enable_encryption: bool = False
    enable_compression: bool = False
    retry_policy: RetryPolicy = RetryPolicy()
    circuit_breaker: CircuitBreakerConfig = CircuitBreakerConfig()
\end{lstlisting}

\subsection{Persistence and Event Sourcing}

When persistence is enabled, every event is written to an aiosqlite database
before being placed in the queue.  This ensures that no event is ever lost,
even in the event of a crash.  A separate \emph{event sourcing log} records
each state mutation as an append-only fact, enabling full audit trails and
deterministic replay.  Replay is implemented as a generator that streams
events from the database in chronological order.

\subsection{Fault Tolerance}

Each subscriber is wrapped with a \emph{retry policy} (exponential backoff)
and a \emph{circuit breaker}.  The circuit breaker transitions through
\texttt{CLOSED}, \texttt{OPEN}, and \texttt{HALF\_OPEN} states based on
failure counts and a configurable timeout (Figure~\ref{fig:circuit}).

\begin{figure}[htbp]
\centering
\begin{tikzpicture}[
    state/.style={draw, rounded corners, minimum width=2cm, minimum height=0.8cm, align=center, font=\small},
    arrow/.style={-{Stealth[length=4mm]}, thick}
]
\node[state, fill=green!20] (closed) {CLOSED};
\node[state, fill=red!20, right=of closed] (open) {OPEN};
\node[state, fill=yellow!20, below=of open] (half) {HALF\_OPEN};
\draw[arrow] (closed) -- node[above] {fail count $\ge$ threshold} (open);
\draw[arrow] (open) -- node[right] {timeout} (half);
\draw[arrow] (half) -- node[left] {success} (closed);
\draw[arrow] (half) -- node[above, sloped] {failure} (open);
\end{tikzpicture}
\caption{Circuit breaker state machine.}
\label{fig:circuit}
\end{figure}

Failed events are moved to a \emph{dead-letter queue} after all retries are
exhausted.  An operator can inspect the dead-letter queue via
\texttt{bus.get\_dead\_letter()} and retry them with
\texttt{bus.retry\_dead\_letter()}.

\subsection{Scalability and Large Payloads}

The event bus is optimised for high‑throughput control messages (typical
payload $< 1$ kB).  For large binary objects such as tensor data, the
recommended pattern is \emph{pass‑by‑reference}: events carry a file
path, shared‑memory handle, or object store identifier, while the bulk
data is transferred out‑of‑band.  This keeps the bus responsive and
avoids serialisation overhead.

When CPU‑bound subscribers are required, the bus can be deployed in a
multi‑process configuration using Redis or Kafka backends, effectively
bypassing the Global Interpreter Lock (GIL).  The in‑memory backend
remains appropriate for single‑agent deployments where all components
reside within the same process.

\subsection{Prometheus Metrics}

An optional Prometheus exporter is built in.  When enabled, it exposes
counters for total events, failed deliveries, and a histogram of
processing latency, allowing integration with standard monitoring stacks.